% ============================================================
% OP_GUT2_repclass_recon_v1.tex  (PROPOSAL-ONLY RECONNAISSANCE)
% OP-GUT2 matter representations: map, tractable kernel, risks.
% NO preprint edits, NO "representations derived" claims.
% ============================================================
\section*{OP-GUT2 representation-classification reconnaissance
(proposal-only)}

\subsection*{0. Scope choice}
The full OP-GUT2 (``derive fermion representations from condensate
dynamics'') is OP3.1-hard and is not attacked.
The tractable kernel, in the AN3/$N_c$-classification mould, is a
\emph{conditional classification}: within a sharply defined
admissible matter class generated by the structures ECT already
postulates (P7 rank-3 colour module, condensate $SU(2)$ doublet
structure, compact-$U(1)$ charge lattice), classify the minimal
anomaly-free chiral completions and compare with the observed
one-generation Standard-Model content.
Outcome shape: ``one SM generation is the/a minimal anomaly-free
chiral set \emph{within the stated class and assumptions}'' ---
a constraint/classification statement, not a derivation of the
representations from dynamics.

\subsection*{1. Existing anchors in the text}
B2$'$ anomaly arithmetic for \eqref{eq:hypercharges}; AN3 solution
structure and the $B-L$ flat direction; the $N_c$-classification
(SM-like family exists for all $N_c\ge2$, Witten parity selects
odd $N_c$; no-overclaim lemma); the charge-lattice theorem
(Level~A) for compact $U(1)$; the Witten even-doublet count under
P7 (twelve $SU(2)$ doublets for $N_{\rm gen}=3$); the zero-mode
programme's rank$\to$charge-3 honesty clause; C2 scope statement
(full anomaly closure conditional on OP-GUT2 $\times$ OP-GUT3b);
the closure-table line ``admissible matter class narrowed by
$SU(3)^3$ and mixed-anomaly cancellation requirements''.

\subsection*{2. Kernel A: module-generated admissible matter
class (definition; bookkeeping level)}
Define the admissible one-generation fermion class as left-handed
Weyl fields valued in tensor products of the postulated modules:
colour $\mathbf{3}$ or $\bar{\mathbf 3}$ (P7 rank-3 module $E$ or
$E^*$), weak $\mathbf{2}$ or $\mathbf 1$ (condensate doublet
structure), and a charge-lattice $U(1)$ label; no higher colour
tensors ($\mathbf 6,\mathbf 8,\dots$) and no higher weak isospin,
\emph{by definition of the class} (these would require module
structures ECT has not postulated), with at most one copy of each
chiral type before counting generations.
This is a definition with a stated rationale, not a theorem; the
exclusion of exotics is a structural input, not an anomaly
consequence.

\subsection*{3. Kernel B: conditional minimal-completion
classification (candidate Level~B after derivation)}
Within the class of Kernel~A impose:
(i) the local anomaly conditions $[SU(3)]^3$, $[SU(3)]^2U(1)$,
$[SU(2)]^2U(1)$, $[U(1)]^3$, and the mixed
$U(1)$--gravitational condition \eqref{eq:mixed_grav_anomaly};
(ii) the Witten $SU(2)$ parity (even doublet count);
(iii) charge-lattice quantisation.
Tasks (own algebra, to be written out in v2):
(a) show a one-generation SM set
$Q\oplus u^c\oplus d^c\oplus L\oplus e^c$ (with or without
$\nu^c$) solves all conditions, reusing AN3/B2$'$;
(b) characterise the solution space: which non-SM chiral
combinations inside the class also pass (expected: essentially
none beyond vector-like pairs and the known $B-L$/hypercharge
normalisation freedom; to be proven, not asserted);
(c) state precisely what additional input beyond anomalies is used
(class definition, minimality, chirality), in a no-overclaim lemma
mirroring the $N_c$ round.
Literature anchors for the known SM-uniqueness folklore, to be
source-verified before any citation: Geng--Marshak (PRD 39, 693
(1989)); Minahan--Ramond--Warner (PRD 41, 715 (1990)); possibly a
modern treatment of anomaly-free chiral sets.
If verification shows their assumptions differ from ours, the
comparison is reported honestly rather than borrowed.

\subsection*{4. What is NOT proposed}
No derivation of representations from condensate dynamics
(OP-GUT2 stays Open; at most ``sharpened'');
no generation counting (OP-generations untouched);
no Yukawa/mass structure; no change to OP-GUT3b status
(hypercharge normalisation freedom remains);
no claim that anomalies alone select the SM content.

\subsection*{5. Risk register}
Forbidden: ``ECT derives the fermion representations'';
``anomalies select the Standard-Model matter content'';
``the SM generation is the unique anomaly-free set'' without the
class/minimality qualifiers.
Safe: ``within the module-generated class and under minimality,
the one-generation SM content is a/the minimal anomaly-free chiral
completion''; ``classification conditional on P7''.
Main hazards: (i) hidden exotic solutions inside the class
(must actually solve the equations, not assume uniqueness);
(ii) borrowing literature uniqueness theorems whose assumptions
(e.g.\ mass terms after EWSB, no gravitational condition) differ
from ours; (iii) silently upgrading OP-GUT2.

\subsection*{6. Questions for GPT}
(Q1)~Approve the two-kernel structure (class definition +
conditional classification), OP-GUT2 status target ``Open;
sharpened'' only?
(Q2)~Class definition: is the no-exotics rationale (no postulated
module structures beyond $E$, doublet, lattice) acceptable as a
stated structural input?
(Q3)~Include $\nu^c$ (lattice-neutral singlet) in the baseline set
or treat as optional branch?
(Q4)~Literature: verify-then-cite Geng--Marshak and
Minahan--Ramond--Warner, or keep the round citation-free and
self-contained?
(Q5)~Proceed as: recon $\to$ your review $\to$ v2 with the full
solution-space algebra $\to$ exact insertion blocks $\to$
implementation?
